% Compile with XeLaTeX: xelatex main.tex
\documentclass[10pt,a4paper]{article}

\usepackage[margin=0.78in,top=0.68in,bottom=0.68in]{geometry}
\usepackage{fontspec}
\usepackage{hyperref}
\usepackage{enumitem}
\usepackage{titlesec}

\setmainfont{TeX Gyre Pagella}
\newfontfamily\chinesefont{Source Han Serif CN}
\hypersetup{hidelinks}
\pagestyle{empty}
\setlength{\parindent}{0pt}
\setlength{\parskip}{0pt}
\setlength{\emergencystretch}{2em}
\linespread{1.03}
\setlist[itemize]{leftmargin=1.2em,topsep=3pt,itemsep=2pt,parsep=0pt,partopsep=0pt}
\titleformat{\section}{\normalsize\bfseries\scshape}{}{0pt}{}[\vspace{-2pt}\titlerule]
\titlespacing*{\section}{0pt}{12pt}{6pt}

\begin{document}
\raggedright

\begin{center}
  {\LARGE\bfseries Jiheng Zhu}\quad{\chinesefont\LARGE 朱霁恒}\\[6pt]
  \href{tel:+85363759311}{+853 63759311}\quad\textbar\quad
  \href{tel:+8618762460011}{+86 18762460011}\\[2pt]
  \href{mailto:sevenszhujiheng@outlook.com}{sevenszhujiheng@outlook.com}\quad\textbar\quad
  \href{https://github.com/sevenszhu}{github.com/sevenszhu}\\[3pt]
  Electronic and Information Engineering Undergraduate
\end{center}

\section*{Education}
\textbf{Macau University of Science and Technology}, Macao, China\hfill Sep 2024 -- Expected Jun 2028\\
Faculty of Innovation Engineering \textbar\ BSc (Electronic and Information Engineering)

\section*{Technical Skills}
\begin{itemize}
  \item \textbf{Programming:} C and C++; able to read Python code.
  \item \textbf{Systems and tools:} Unix-like systems (Linux and macOS), common shell commands, and Git version control.
  \item \textbf{PCB and circuit design:} Schematic capture and two-layer PCB layout in EasyEDA; experience with Altium Designer, datasheet interpretation, component selection, and circuit analysis.
  \item \textbf{Testing and prototyping:} Circuit measurement and troubleshooting with a multimeter and oscilloscope; PCB soldering, testing, and modification.
  \item \textbf{Embedded systems:} Hands-on STM32 circuit work and PWM-driven Boost converter prototyping.
\end{itemize}

\section*{Engineering Experience}
\textbf{Wuxi Huapu Microelectronics Co., Ltd.}, Wuxi, China\hfill Aug 2026\\
\textit{Intern}
\begin{itemize}
  \item Reviewed and updated bills of materials (BOMs), cross-checking component identifiers, specifications, and quantities and recording changes for engineering review.
  \item Supported PCB reverse engineering of reference boards by identifying components, tracing circuit connections, and documenting placement and connectivity for schematic and layout reconstruction.
\end{itemize}

\textbf{Power Electronics Laboratory}, Macau University of Science and Technology\hfill Jun -- Jul 2026\\
\textit{Undergraduate Student}
\begin{itemize}
  \item Built and iteratively tested a microcontroller-controlled Boost converter prototype, assembling control and power-stage circuits and generating PWM signals to drive MOSFET switching.
  \item Verified stable PWM output, investigated power-stage switching waveforms and boost output, and documented tests and troubleshooting steps in a laboratory report.
\end{itemize}

\section*{Research Experience}
\textbf{Nanomaterials Laboratory}, Macau University of Science and Technology\hfill May -- Oct 2025\\
\textit{Undergraduate Student}
\begin{itemize}
  \item Assisted with nanomaterials and electrocatalysis experiments, carrying out laboratory procedures, recording experimental conditions, and organizing results.
  \item Organized and analyzed experimental data, creating visualizations to support interpretation of results.
\end{itemize}

\section*{Publications}
Z. Qu, S. Zhang, \textbf{J. Zhu}, et al. ``Construction of stable Cu--O--Fe bridging structure for highly efficient electrocatalytic nitrate reduction to ammonia.'' \textit{Chemical Communications}, 2026. DOI: \href{https://doi.org/10.1039/D6CC01243D}{10.1039/D6CC01243D}.\\[3pt]
T. Hu, H. Yang, Z. Li, \textbf{J. Zhu}, et al. ``Interfacial Photophysics and Solvent-triggered Phase Reconstruction in Cs$_4$PbBr$_6$ Quantum Dots for Optical/Electrical Dual-Mode Sensing.'' \textit{Small}, 2026. DOI: \href{https://doi.org/10.1002/smll.75924}{10.1002/smll.75924}.

\section*{Relevant Coursework}
Digital Circuits, Analog Circuits, Circuit Analysis, Computer Organization, FPGA System Design, and Signals and Systems.

\section*{Professional Strengths}
Collaborative communicator who coordinates effectively with others, welcomes feedback, and follows through on improvements.

\end{document}
